% ---------------------------------------------------------------------------
% Study report · NeuralProphet inclination prediction
%
% Build: pdflatex neuralprophet_inclination_prediction_report.tex (twice)
% Figures and numerical results are produced by the paired notebook.
% ---------------------------------------------------------------------------
\documentclass[11pt,a4paper]{article}

\newcommand{\runninghead}{NeuralProphet prediction of inclination change}

% ---------------------------------------------------------------------------
% reportstyle.tex — shared preamble for the study reports under studies/.
%
% Kept deliberately inside the package set shipped with TeX Live "basic":
% no tcolorbox, no siunitx, no enumitem, no titlesec. Everything below
% compiles with a stock pdflatex installation.
%
% Usage, from a report directory one level below a study folder:
%     % ---------------------------------------------------------------------------
% reportstyle.tex — shared preamble for the study reports under studies/.
%
% Kept deliberately inside the package set shipped with TeX Live "basic":
% no tcolorbox, no siunitx, no enumitem, no titlesec. Everything below
% compiles with a stock pdflatex installation.
%
% Usage, from a report directory one level below a study folder:
%     % ---------------------------------------------------------------------------
% reportstyle.tex — shared preamble for the study reports under studies/.
%
% Kept deliberately inside the package set shipped with TeX Live "basic":
% no tcolorbox, no siunitx, no enumitem, no titlesec. Everything below
% compiles with a stock pdflatex installation.
%
% Usage, from a report directory one level below a study folder:
%     \input{../../_shared/reportstyle.tex}
%     \graphicspath{{../outputs/}}
% ---------------------------------------------------------------------------

\usepackage[T1]{fontenc}
\usepackage[utf8]{inputenc}
\usepackage{textcomp}
\usepackage{lmodern}
\usepackage{microtype}
\usepackage[margin=2.4cm,headheight=15pt]{geometry}
\usepackage{graphicx}
\usepackage{booktabs}
\usepackage{tabularx}
\usepackage{array}
\usepackage{amsmath}
\usepackage{xcolor}
\usepackage[font=small,labelfont=bf,justification=justified]{caption}
\usepackage{float}
\usepackage{fancyhdr}
\usepackage[hidelinks]{hyperref}

% --- Palette ---------------------------------------------------------------
\definecolor{rulegrey}{HTML}{B9C0C7}
\definecolor{fillgrey}{HTML}{F2F4F6}
\definecolor{failred}{HTML}{A5202B}
\definecolor{passgreen}{HTML}{1B6B3A}
\definecolor{accent}{HTML}{1F4E79}

% --- Semantic shorthands ---------------------------------------------------
\newcommand{\degC}{\textdegree C}
\newcommand{\mdegC}{mdeg/\textdegree C}
\newcommand{\fail}{\textcolor{failred}{\textbf{fail}}}
\newcommand{\pass}{\textcolor{passgreen}{\textbf{pass}}}
\newcommand{\worse}{\textcolor{failred}{worse}}
\newcommand{\better}{\textcolor{passgreen}{better}}

% --- Highlighted panel -----------------------------------------------------
% #1 = heading, #2 = body
\newcommand{\panel}[2]{%
  \begin{center}%
  \setlength{\fboxsep}{9pt}%
  \fcolorbox{rulegrey}{fillgrey}{%
    \begin{minipage}{0.93\textwidth}%
      {\bfseries\color{accent}#1}\par\medskip
      #2
    \end{minipage}}%
  \end{center}}

% --- Numeric column: right aligned, fixed width ----------------------------
\newcolumntype{R}[1]{>{\raggedleft\arraybackslash}p{#1}}
\newcolumntype{L}[1]{>{\raggedright\arraybackslash}p{#1}}
\newcolumntype{Y}{>{\raggedright\arraybackslash}X}

% --- Table look ------------------------------------------------------------
\renewcommand{\arraystretch}{1.15}
\setlength{\tabcolsep}{6pt}

% --- Headers ---------------------------------------------------------------
\pagestyle{fancy}
\fancyhf{}
\fancyhead[L]{\small\itshape\runninghead}
\fancyhead[R]{\small\thepage}
\renewcommand{\headrulewidth}{0.4pt}

% --- Section spacing without titlesec --------------------------------------
\makeatletter
\renewcommand\section{\@startsection{section}{1}{\z@}%
  {-3.2ex \@plus -1ex \@minus -.2ex}{2.0ex \@plus.2ex}%
  {\normalfont\Large\bfseries\color{accent}}}
\renewcommand\subsection{\@startsection{subsection}{2}{\z@}%
  {-2.6ex \@plus -1ex \@minus -.2ex}{1.4ex \@plus.2ex}%
  {\normalfont\large\bfseries}}
\makeatother

% --- Figure defaults -------------------------------------------------------
\setkeys{Gin}{width=\linewidth}
\captionsetup[figure]{skip=6pt}

%     \graphicspath{{../outputs/}}
% ---------------------------------------------------------------------------

\usepackage[T1]{fontenc}
\usepackage[utf8]{inputenc}
\usepackage{textcomp}
\usepackage{lmodern}
\usepackage{microtype}
\usepackage[margin=2.4cm,headheight=15pt]{geometry}
\usepackage{graphicx}
\usepackage{booktabs}
\usepackage{tabularx}
\usepackage{array}
\usepackage{amsmath}
\usepackage{xcolor}
\usepackage[font=small,labelfont=bf,justification=justified]{caption}
\usepackage{float}
\usepackage{fancyhdr}
\usepackage[hidelinks]{hyperref}

% --- Palette ---------------------------------------------------------------
\definecolor{rulegrey}{HTML}{B9C0C7}
\definecolor{fillgrey}{HTML}{F2F4F6}
\definecolor{failred}{HTML}{A5202B}
\definecolor{passgreen}{HTML}{1B6B3A}
\definecolor{accent}{HTML}{1F4E79}

% --- Semantic shorthands ---------------------------------------------------
\newcommand{\degC}{\textdegree C}
\newcommand{\mdegC}{mdeg/\textdegree C}
\newcommand{\fail}{\textcolor{failred}{\textbf{fail}}}
\newcommand{\pass}{\textcolor{passgreen}{\textbf{pass}}}
\newcommand{\worse}{\textcolor{failred}{worse}}
\newcommand{\better}{\textcolor{passgreen}{better}}

% --- Highlighted panel -----------------------------------------------------
% #1 = heading, #2 = body
\newcommand{\panel}[2]{%
  \begin{center}%
  \setlength{\fboxsep}{9pt}%
  \fcolorbox{rulegrey}{fillgrey}{%
    \begin{minipage}{0.93\textwidth}%
      {\bfseries\color{accent}#1}\par\medskip
      #2
    \end{minipage}}%
  \end{center}}

% --- Numeric column: right aligned, fixed width ----------------------------
\newcolumntype{R}[1]{>{\raggedleft\arraybackslash}p{#1}}
\newcolumntype{L}[1]{>{\raggedright\arraybackslash}p{#1}}
\newcolumntype{Y}{>{\raggedright\arraybackslash}X}

% --- Table look ------------------------------------------------------------
\renewcommand{\arraystretch}{1.15}
\setlength{\tabcolsep}{6pt}

% --- Headers ---------------------------------------------------------------
\pagestyle{fancy}
\fancyhf{}
\fancyhead[L]{\small\itshape\runninghead}
\fancyhead[R]{\small\thepage}
\renewcommand{\headrulewidth}{0.4pt}

% --- Section spacing without titlesec --------------------------------------
\makeatletter
\renewcommand\section{\@startsection{section}{1}{\z@}%
  {-3.2ex \@plus -1ex \@minus -.2ex}{2.0ex \@plus.2ex}%
  {\normalfont\Large\bfseries\color{accent}}}
\renewcommand\subsection{\@startsection{subsection}{2}{\z@}%
  {-2.6ex \@plus -1ex \@minus -.2ex}{1.4ex \@plus.2ex}%
  {\normalfont\large\bfseries}}
\makeatother

% --- Figure defaults -------------------------------------------------------
\setkeys{Gin}{width=\linewidth}
\captionsetup[figure]{skip=6pt}

%     \graphicspath{{../outputs/}}
% ---------------------------------------------------------------------------

\usepackage[T1]{fontenc}
\usepackage[utf8]{inputenc}
\usepackage{textcomp}
\usepackage{lmodern}
\usepackage{microtype}
\usepackage[margin=2.4cm,headheight=15pt]{geometry}
\usepackage{graphicx}
\usepackage{booktabs}
\usepackage{tabularx}
\usepackage{array}
\usepackage{amsmath}
\usepackage{xcolor}
\usepackage[font=small,labelfont=bf,justification=justified]{caption}
\usepackage{float}
\usepackage{fancyhdr}
\usepackage[hidelinks]{hyperref}

% --- Palette ---------------------------------------------------------------
\definecolor{rulegrey}{HTML}{B9C0C7}
\definecolor{fillgrey}{HTML}{F2F4F6}
\definecolor{failred}{HTML}{A5202B}
\definecolor{passgreen}{HTML}{1B6B3A}
\definecolor{accent}{HTML}{1F4E79}

% --- Semantic shorthands ---------------------------------------------------
\newcommand{\degC}{\textdegree C}
\newcommand{\mdegC}{mdeg/\textdegree C}
\newcommand{\fail}{\textcolor{failred}{\textbf{fail}}}
\newcommand{\pass}{\textcolor{passgreen}{\textbf{pass}}}
\newcommand{\worse}{\textcolor{failred}{worse}}
\newcommand{\better}{\textcolor{passgreen}{better}}

% --- Highlighted panel -----------------------------------------------------
% #1 = heading, #2 = body
\newcommand{\panel}[2]{%
  \begin{center}%
  \setlength{\fboxsep}{9pt}%
  \fcolorbox{rulegrey}{fillgrey}{%
    \begin{minipage}{0.93\textwidth}%
      {\bfseries\color{accent}#1}\par\medskip
      #2
    \end{minipage}}%
  \end{center}}

% --- Numeric column: right aligned, fixed width ----------------------------
\newcolumntype{R}[1]{>{\raggedleft\arraybackslash}p{#1}}
\newcolumntype{L}[1]{>{\raggedright\arraybackslash}p{#1}}
\newcolumntype{Y}{>{\raggedright\arraybackslash}X}

% --- Table look ------------------------------------------------------------
\renewcommand{\arraystretch}{1.15}
\setlength{\tabcolsep}{6pt}

% --- Headers ---------------------------------------------------------------
\pagestyle{fancy}
\fancyhf{}
\fancyhead[L]{\small\itshape\runninghead}
\fancyhead[R]{\small\thepage}
\renewcommand{\headrulewidth}{0.4pt}

% --- Section spacing without titlesec --------------------------------------
\makeatletter
\renewcommand\section{\@startsection{section}{1}{\z@}%
  {-3.2ex \@plus -1ex \@minus -.2ex}{2.0ex \@plus.2ex}%
  {\normalfont\Large\bfseries\color{accent}}}
\renewcommand\subsection{\@startsection{subsection}{2}{\z@}%
  {-2.6ex \@plus -1ex \@minus -.2ex}{1.4ex \@plus.2ex}%
  {\normalfont\large\bfseries}}
\makeatother

% --- Figure defaults -------------------------------------------------------
\setkeys{Gin}{width=\linewidth}
\captionsetup[figure]{skip=6pt}

\graphicspath{{../outputs/}}

\title{\vspace{-1.2cm}\textbf{Can the on-structure sensors predict inclination?}\\[2mm]
\large A chronological NeuralProphet experiment on same-time estimation,
one-week forecasting and missing-hour reconstruction}
\author{Study report · \texttt{studies/04\_neuralprophet\_inclination\_prediction/}}
\date{}

\begin{document}
\maketitle
\thispagestyle{fancy}



\tableofcontents



% ===========================================================================
\section{Introduction}

Study~01 read the raw archive, decided which of the recorded values are
measurements, and exported the verdict-aware archive that every later study
reads. Its final inclination product --- compensated for temperature, anchored
once across the whole record, cleaned of impulsive noise, and carrying a flag
that marks every value it replaced by interpolation --- is the only inclination
series used here. The intermediate versions of that channel are not used, and no
value the flag marks as an interpolation is treated as a measurement. This study
inherits that verdict rather than revisiting it.

The subject is one sensor at one place: the inclinometer at station~02, on the
medieval wall at Porta di Sant'Ubaldo. The question is operational rather than
diagnostic. Given the record as it actually stands, can the next value of the
inclination be predicted --- from the inclination's own history, from the other
variables the instrument package on the structure measures beside it, or from
external environmental proxies?

The record is not taken whole. The archive contains a 271-day outage that begins
on 23~September 2022 and whose last missing day is 20~June 2023, across which
nothing at all about the structure's behaviour was recorded. This study works
only on what follows: the stretch running from the resumption of recording in
June~2023 to the end of the archive. The choice has two reasons. That stretch is
the part of the record closest to the present, and therefore the part a
monitoring system deployed today would actually have to work with. It is also
the only part over which the complete instrument package is available, since the
wall-temperature probe and the pyranometer were installed on 21~February 2025
and record nothing before that date.

Choosing that window does not buy a continuous series. Four further whole-day
outages fall inside it --- 40~days from 5~May 2024, 42~days from 28~August 2024,
17~days from 26~September 2025, and 103~days from 7~March 2026 --- and beneath
them lies a population of shorter dropouts whose size is not yet known. Over the
27,769 hourly slots the window spans, the inclinometer returns an accepted value
in 76.2\% of them, and air temperature and relative humidity in 78.4\%. The
wall-temperature probe and the pyranometer reach only 17.7\% and
18.9\%: they exist for less than half the window to begin with, and they are
interrupted inside the part where they do exist. What none of those numbers says
is how the missing time is \emph{distributed}: whether it arrives as isolated
absent samples, as hours, or as days. Study~01 recorded the omission explicitly,
noting that the anatomy of the gaps is a natural extension of it and was not
attempted in it. Establishing that anatomy is the first task of this study,
because the shape of the missing data decides what kind of statistical problem
filling it is: recovering a scattered sample from its immediate neighbours is a
different proposition from reconstructing a season.

The study therefore proceeds in four movements. It first looks at the window
whole, plotting every variable the on-structure package records over it on one
shared clock, with the gaps left as gaps and nothing interpolated before
drawing. It then inventories those gaps, classifying the missing time by
duration so that the scattered absences are separated from the outages. It then
weighs the ways each class might be filled, and what it would cost to be wrong.
Only then does it put the prediction question: what the next inclination value
can be inferred from, using the on-structure measurements.

\begin{figure}[H]
\centering
\includegraphics{NP_F01_on_structure_record.png}
\caption{Every variable the instrument package on the wall records, over the
window this study works on, on one shared clock. Nothing is interpolated and
nothing is dropped before drawing: a gap in the record is a gap in the figure.
Inclination, air temperature and relative humidity run from the resumption of
recording in June~2023; wall temperature and solar radiation begin only with the
package installed on 21~February 2025, and are themselves interrupted after it.
The package also records its own supply voltage, which is instrument
housekeeping rather than a measurement of the wall or of its environment, and is
not drawn here.}
\label{fig:on-structure-record}
\end{figure}

% ===========================================================================
\section{The record this study works on}
\label{sec:record}

The window runs from 21~June 2023 to 20~August 2026 and spans 83,305 slots on
the twenty-minute grid at which the acquisition system writes. That grid, rather
than an hourly aggregate, is the one every result below is stated on, because it
is the rate at which a deployed system receives a new reading and therefore the
rate at which it must decide whether that reading is the expected one.

Table~\ref{tab:coverage} gives what each channel returns over the window.
The inclinometer delivers an accepted value in 75.5\% of the slots, air
temperature and relative humidity in 77.8\% of them. Wall temperature and solar
radiation reach 17.6\% and 18.8\%: the instruments that record them were
installed on 21~February 2025 and so do not exist for the first three fifths of
the window at all. That is the whole of the reason those two channels are
excluded from every model in this study. Requiring them would not cost a
percentage point of coverage but three fifths of the record, and every result
would then have to be reported twice, on two windows that cannot be compared.
They remain in Figure~\ref{fig:on-structure-record} because what the package
records is part of what this study has to describe, and a reader is owed the
sight of the channels it set aside.

\begin{table}[H]
\centering
\caption{Accepted values and coverage per channel over the window, on the
twenty-minute grid. Read from \texttt{NP\_01\_window\_coverage.csv}.}
\label{tab:coverage}
\begin{tabular}{lrr}
\toprule
Channel & Accepted slots & Coverage \\
\midrule
\input{../outputs/NP_01_body.tex}\\
\bottomrule
\end{tabular}
\end{table}

Air temperature and relative humidity cost this study nothing. Their
missingness is nested inside the inclinometer's: every slot in which the
inclination is accepted is a slot in which both are accepted too. Requiring all
three therefore leaves exactly the same rows, in exactly the same contiguous
runs, as requiring the inclination alone. That is a property of this record and
not a general one, and it is the reason the predictor set could be chosen on
physical grounds rather than by trading coverage against information.

% ===========================================================================
\section{The anatomy of what is missing}
\label{sec:gaps}

A coverage percentage says how much time is missing. It does not say how that
time is shaped, and the shape is what decides whether filling it is
interpolation or reconstruction. Over the window the inclination record is
interrupted 1,647 times, for a total of 6,791 hours.

Those interruptions are distributed with extreme asymmetry, and in two
directions at once. Four gaps longer than seven days --- 40~days from 5~May
2024, 42~days from 28~August 2024, 17~days from 26~September 2025 and 103~days
from 7~March 2026 --- carry 72.8\% of all the missing time between them, while
being 0.2\% of the gaps. At the other end, 1,280 gaps of an hour or less carry
only 8.6\% of the missing time while being 77.7\% of the gaps.
Table~\ref{tab:gapclasses} gives the full classification.

\begin{table}[H]
\centering
\caption{The missing time by gap duration: how many gaps of each class, and how
many hours they carry between them. Read from
\texttt{NP\_02\_gap\_inventory.csv}.}
\label{tab:gapclasses}
\begin{tabular}{lrr}
\toprule
Gap class & Gaps & Hours \\
\midrule
\input{../outputs/NP_02_body.tex}\\
\bottomrule
\end{tabular}
\end{table}

The second number matters more than the first, and it is the one a coverage
figure hides. The four long outages are conspicuous, and no honest method would
try to reconstruct 103~days of a structure's behaviour from its neighbours. It
is the short dropouts that determine what can be modelled, because they are what
breaks \emph{contiguity}, and contiguity is the currency an autoregressive
window spends. A model asked for 24~hours of history before each prediction
cannot use a run of data shorter than that, however complete the run is.

Table~\ref{tab:survival} prices that directly. The record divides into 1,647
contiguous runs whose median length is 1.7~hours. Requiring 24~hours of history
and forecasting 24~hours ahead leaves 101 of those runs usable and 15,000
prediction windows; requiring 48~hours of history for the same forecast leaves
58 runs and 9,092 windows. The autoregressive window used in
Section~\ref{sec:forecast} is chosen from this table rather than from habit.

\begin{table}[H]
\centering
\caption{What each autoregressive window costs in usable data. A contiguous run
survives only if it is longer than the lag plus the forecast it must support.
Read from \texttt{NP\_03\_segment\_survival.csv}.}
\label{tab:survival}
\begin{tabular}{rrrrr}
\toprule
Lag (h) & Forecast (h) & Runs & Surviving & Windows \\
\midrule
\input{../outputs/NP_03_body.tex}\\
\bottomrule
\end{tabular}
\end{table}

\begin{figure}[H]
\centering
\includegraphics[width=\textwidth]{NP_F02_gap_anatomy.png}
\caption{The anatomy of the missing time. The gaps are counted by duration class
and weighted by the time each class carries, which is what separates the many
short dropouts from the few long outages.}
\label{fig:gap-anatomy}
\end{figure}

\begin{figure}[H]
\centering
\includegraphics[width=\textwidth]{NP_F03_segment_survival.png}
\caption{Contiguous runs surviving each autoregressive window, and the
prediction windows they yield. The cost of a longer memory on a fragmented
record is not linear: it falls on the number of runs that qualify at all.}
\label{fig:segment-survival}
\end{figure}

% ===========================================================================
\section{Method}
\label{sec:method}

Each subsection below states one decision, the evidence that forced it, and what
would have gone wrong under the alternative. They are given before any result
that depends on them.

\subsection{Why the level is decomposed but never scored}

The inclination level has a lag-one autocorrelation of 0.9977 at twenty minutes.
Any error metric computed against it is therefore a measurement of the sampling
interval rather than of a model: predicting that the next value equals the
current one is almost exactly right, and a model that beats that baseline by a
few per cent has demonstrated nothing about the structure. The level is
decomposed, because a decomposition of the level is what says what the record is
\emph{made of}; it is never scored as a forecast.

\subsection{Why forecasting is done on the hourly change}

Differencing removes the autocorrelation but introduces a different problem, and
the cadence decides which of the two dominates. Table~\ref{tab:cadence} gives
both. At twenty minutes the first difference has a lag-one autocorrelation of
$-0.058$, which is the signature of additive measurement noise on the level
rather than of structural memory, and it correlates with the change in air
temperature at $-0.778$. At one hour the same difference retains $+0.349$ of
memory and couples to air temperature at $-0.892$. The hourly change is
therefore the cadence at which the increment carries signal rather than noise,
and it is the target of the forecasting model.

\begin{table}[H]
\centering
\caption{The evidence for the two cadences. Read from
\texttt{NP\_04\_cadence\_evidence.csv}.}
\label{tab:cadence}
\begin{tabular}{lrrrrr}
\toprule
Cadence & Level $r_1$ & Change $r_1$ & Change s.d. & $\rho(\Delta$inc$,\Delta$tair$)$ & Drift (mdeg/yr) \\
\midrule
\input{../outputs/NP_04_body.tex}\\
\bottomrule
\end{tabular}
\end{table}

\begin{figure}[H]
\centering
\includegraphics[width=\textwidth]{NP_F04_cadence_evidence.png}
\caption{Why the forecasting target is the hourly change rather than the level
or the twenty-minute change.}
\label{fig:cadence}
\end{figure}

\subsection{Why the anomaly test is not differenced, and why contemporaneous
regressors are not leakage}

The two questions this study puts to the record are different questions, and
they are entitled to different rules. Asking \emph{how far ahead can this be
forecast} requires that the model see only the past, and the forecasting model of
Section~\ref{sec:forecast} is held to that: it consumes lagged predictor values
only. Asking \emph{is this reading the expected one} does not. A monitoring
system judging a newly arrived inclination has, in the same acquisition record,
the air temperature and relative humidity measured at the same instant --- they
arrive together, twenty minutes at a time. Consuming them is not leakage but a
description of what the system actually holds when it must decide. The two
models are never scored on the same table, and no number from one is quoted
against the other.

\subsection{Why two models, and why the decomposition carries no autoregression}

With autoregressive lags enabled, the autoregressive component absorbs most of
the diurnal structure and the seasonal component becomes whatever periodicity is
left over after it, rather than the wall's thermal cycle. A decomposition meant
to be read as physics must therefore carry no autoregressive term, which is
exactly the term a forecast needs. One model cannot do both jobs, so there are
two.

\subsection{How gaps are honoured}

No value is imputed anywhere in this study. Every contiguous run of accepted
observations carries its own segment identifier, and the models are told about
those identifiers, so that no fitted window and no prediction window spans a
gap. This matters because NeuralProphet's own defaults do not honour it: left to
itself the library fills gaps of at least thirty hours, and a decomposition
fitted over fabricated data attributes to the structure a shape the structure
never produced.

\subsection{Changepoints on covered time, and no era offset}

The trend's changepoints are placed at quantiles of the \emph{observed}
timestamps rather than uniformly along the axis. A changepoint inside a 103-day
outage is constrained by no data and is free to take any value the extrapolation
on either side of it requires. Instrument eras are not modelled at all: Study~01
compensated and anchored the record once across the changeover of
21~February 2025, and this study takes that product as delivered. No era offset
is estimated and none is subtracted.

\subsection{The metrics, and what each is for}

Point accuracy is reported as mean absolute error and root mean square error in
millidegrees, and as MASE against a naive scale computed on training rows only,
so that the scale-free figure is not self-referential. An interval is reported
on three counts rather than one, because coverage alone can always be bought by
widening: how often it contains the observation, how wide it had to be to manage
that, and the Winkler interval score which penalises both together.

% ===========================================================================
\section{What the record is made of}
\label{sec:decomposition}

Table~\ref{tab:shares} gives the variance each fitted component carries and the
peak-to-peak range it spans. The result is dominated by two things: a trend
carrying 48.9\% of the variance and spanning 213~mdeg, and a residual carrying
30.4\%. Air temperature accounts for 12.2\% and the annual cycle 8.4\%. The
daily seasonality carries almost nothing --- 0.03\% --- and that is a
consequence of the model's construction rather than a finding about the wall:
air temperature enters as a contemporaneous regressor, so it absorbs the diurnal
variation, and what the daily seasonal term is left to explain is the small
residue the temperature does not.

\begin{table}[H]
\centering
\caption{Variance share and peak-to-peak range of each fitted component, over
the training window. Read from \texttt{NP\_05\_component\_shares.csv}.}
\label{tab:shares}
\begin{tabular}{lrrr}
\toprule
Component & Variance & Share & Peak-to-peak (mdeg) \\
\midrule
\input{../outputs/NP_05_body.tex}\\
\bottomrule
\end{tabular}
\end{table}

\begin{figure}[H]
\centering
\includegraphics[width=\textwidth]{NP_F05_decomposition_stack.png}
\caption{The decomposition of the compensated, cleaned inclination level into
trend, annual and daily cycles, the response to the measured environment, and
the remainder. Lines break across dropped outages rather than interpolating
through them.}
\label{fig:decomposition}
\end{figure}

\begin{figure}[H]
\centering
\includegraphics[width=\textwidth]{NP_F06_daily_cycle_and_response.png}
\caption{The two interpretable components on their own axes over a fortnight:
the daily cycle the model fitted, and the response it attributes to air
temperature and relative humidity.}
\label{fig:daily-cycle}
\end{figure}

\subsection{Confrontation with Study 03}

The thermal gain this decomposition learns is an internal quantity of a fitted
model, and it is worth only as much as its agreement with an independent
measurement. Study~03 measured the diurnal-band gain of the same channel against
the same air temperature by a lag-and-gain scan, a method sharing no machinery
with this one, and obtained $-2.79$~\mdegC. This decomposition returns
$-2.56$~\mdegC. Table~\ref{tab:gains} places them side by side, together with the
two weaker external proxies Study~03 also scanned.

\begin{table}[H]
\centering
\caption{The learned thermal gain against Study~03's independent measurements.
Read from \texttt{NP\_06\_learned\_gains.csv}.}
\label{tab:gains}
\begin{tabular}{llr}
\toprule
Source & Gain (\mdegC) & Method \\
\midrule
\input{../outputs/NP_06_body.tex}\\
\bottomrule
\end{tabular}
\end{table}

The two agree to within 8\%, by methods with nothing in common, and both carry
the negative sign the embankment geometry predicts. That is the strongest
external check this study has, and it is the reason the decomposition's thermal
term can be read as a description of the wall rather than as a fitting artefact.
It is a statement about the \emph{compensated} channel: what the wall does after
Study~01's correction, which is the only quantity a monitoring system ever sees.
Whether that correction is itself correct is Study~01's question and is not
reopened here.

\subsection{Which components survive re-estimation}

A component estimated once is an assertion. Table~\ref{tab:stability} re-fits the
model on each chronological third of the training period and compares the two
interpretable quantities across them. The thermal gain is stable: $-2.47$,
$-2.33$ and $-2.39$~\mdegC, keeping its sign and its magnitude across three
disjoint stretches. The trend is not: $+38.6$, $+4.8$ and $+45.5$~mdeg/yr, each
third disagreeing with the whole-window estimate of $-2.77$~mdeg/yr in sign as
well as size.

\begin{table}[H]
\centering
\caption{The two interpretable quantities re-estimated on disjoint thirds of the
training period. Read from \texttt{NP\_16\_component\_stability.csv}.}
\label{tab:stability}
\begin{tabular}{lrr}
\toprule
Block & Thermal gain (\mdegC) & Trend (mdeg/yr) \\
\midrule
\input{../outputs/NP_16_body.tex}\\
\bottomrule
\end{tabular}
\end{table}

The contrast is the finding. The thermal gain behaves like a physical constant of
the structure; the trend behaves like an estimate that has not converged. The
drift is therefore quoted once, in-sample, at $-2.77$~mdeg/yr over 21~June 2023
to 1~September 2025, and no extrapolation of it is used anywhere in this study.
The monitoring model of Section~\ref{sec:expectation} carries no trend at all for
the same reason.

% ===========================================================================
\section{Is this reading the one that was expected?}
\label{sec:expectation}

Three expectations were built and scored against each other, and
Table~\ref{tab:nowcast} reports all three rather than the best. Two are frozen:
a model fitted once at the training origin and applied thereafter, in the
attribution configuration that carries a trend and in the monitoring
configuration that does not. The third is refitted every thirty days on all
history available at each origin.

\begin{table}[H]
\centering
\caption{The expectation scored three ways. Coverage and width refer to the
nominal 90\% interval. Read from \texttt{NP\_07\_nowcast\_metrics.csv}.}
\label{tab:nowcast}
\begin{tabular}{lrrrrrrr}
\toprule
Fit & $n$ & MAE & RMSE & Bias & MASE & Coverage & Width \\
\midrule
\input{../outputs/NP_07_body.tex}\\
\bottomrule
\end{tabular}
\end{table}

The refitted expectation wins on every count, and by margins that are not
marginal: mean absolute error 18.2~mdeg against 30.6 for the frozen monitoring
fit and 58.8 for the frozen attribution fit, bias $-10.1$~mdeg against $-28.8$
and $+57.5$, interval score 164 against 511 and 1197. The reason is visible in
the bias column. A fit frozen at the training origin sits nearly 29~mdeg away
from the record over the following year, and 82\% of its mean square error is in
that single offset rather than in scatter --- it is not noisy, it is stale.
Refitting every thirty days holds the expectation to the drift a month can
accumulate, about 0.23~mdeg at the measured rate, and it is also what a deployed
system would actually do.

The interval is empirical rather than the model's own. NeuralProphet's quantile
regression fits the residual's core, and this residual has a tight core with fat
tails; its nominal 90\% interval covers 68.7\% even in-sample. Conformal
quantiles, calibrated on residuals from before 1~June 2025 and never on rows the
interval is scored against, replace it.

They do not repair it. The reported coverage is \textbf{67.7\%} against a
nominal 90\%, at a median width of 77.9~mdeg. This is stated as measured and is
not adjusted. Two things account for it, and both are consequences of decisions
taken for other reasons. The calibration stretch ends at a fixed date while the
minimum training history was raised to a full annual cycle, which shortened that
stretch from roughly seventeen months to eleven. And the monitoring model
carries no trend by design, so the wall's own slow movement remains in the
residual and continues after the calibration stretch ends --- the interval is
calibrated on a quieter record than the one it is then asked to cover. An
interval that undercovers by this much is a statement about this expectation,
and Section~\ref{sec:limitations} records it as one.

\begin{figure}[H]
\centering
\includegraphics[width=\textwidth]{NP_F07_observed_vs_expected.png}
\caption{One month of the observed inclination against what the measured
environment predicted, with the conformal interval around the expectation.}
\label{fig:observed-expected}
\end{figure}

% ===========================================================================
\section{Judging a departure}
\label{sec:monitoring}

\subsection{The reference window, and what had to be excluded from it}

Every control limit below is a multiple of a centre and a scale estimated on a
reference window, so a window containing a departure would calibrate the detector
against the very thing it is meant to find. Two stretches are therefore excluded,
for two different reasons.

The window ends on 1~June 2025, before the stretch Study~01 flagged in summer
2026. That is a precaution about calibration and explicitly not a test: this
study makes no claim, anywhere, about whether this detector finds that stretch.

The window also begins later than the monitored record could allow, on
1~October 2024, and this exclusion was forced by measurement rather than
anticipated. The residual's thirty-day rolling mean runs $-60.5$, $-61.0$,
$-34.2$ and $-35.5$~mdeg over June to September 2024 and then settles to
$+3.4$~mdeg and stays within roughly $\pm20$~mdeg for the remainder of the
record. The cause is identifiable: those early refits train on barely one
calendar annual cycle, and one calendar year of this record contains outages of
40, 42 and 17~days, so their annual term is fitted to substantially less than one
cycle and is not identifiable. Those months are excluded from the monitored
record as well as from the reference window. An expectation wrong by 60~mdeg for
a reason internal to the model is not a structural departure, and reporting it as
one would be a false claim.

\subsection{The limit width, and why it is far above the textbook range}

The limit is not taken from convention. A false-alarm budget is fixed first ---
ninety days of watched time per false alarm --- and the limit width is then swept
until one meets it, so that every detection figure in this study is quoted at a
stated budget rather than at whatever rate a conventional setting happens to
give.

The width that meets it is 14.75 standard deviations, which is far outside
anything a control-chart text would offer, and the reason is a property of this
residual that must be stated plainly. EWMA and CUSUM limits are derived for
independent observations. This residual has a lag-one autocorrelation of
\textbf{0.995 at twenty minutes} and still \textbf{0.974 at twenty-four hours}.
Against a series that correlated, a textbook width between two and five standard
deviations does not buy months between false alarms but days: the best any width
in that range achieves here is 15.7~days against the ninety asked for. Widening
the swept range is not the same as widening the limit until the alarms stop ---
the budget was fixed before the sweep, and the width is whatever meets it --- but
the resulting number should be read as a property of the residual rather than as
a conventional control limit.

At 14.75 the achieved run length is 122~days per false alarm. It rests on two
episodes over 244 watched days, and with a count that small it bounds the
false-alarm rate rather than measuring it. It is quoted here to one significant
figure of confidence and no more.

\subsection{What the detector then reported}

Over the monitored record the joint chart raises 14 alarming episodes covering
350 of 678 watched days. Table~\ref{tab:episodes} gives them. One episode
dominates: 288 consecutive days from 24~September 2025, at a mean standardised
departure of $+5.2$.

\begin{table}[H]
\centering
\caption{Alarming episodes over the monitored record, where the EWMA and CUSUM
charts agree within six hours. Read from
\texttt{NP\_09\_alarm\_episodes.csv}.}
\label{tab:episodes}
\begin{tabular}{llrrr}
\toprule
Start & End & Duration (h) & Mean $z$ & Peak $|z|$ \\
\midrule
\input{../outputs/NP_09_body.tex}\\
\bottomrule
\end{tabular}
\end{table}

A detector that spends half its watched time in alarm is reporting something
about the model as much as about the wall. The monitoring fit carries no trend
by deliberate choice, so the wall's slow movement is left in the residual where a
monitor can see it; the consequence is that when the wall does move slowly and
persistently, the chart alarms and keeps alarming. That is the intended
behaviour, and it is also why the episode table cannot be read as a list of
distinct structural events.

\begin{figure}[H]
\centering
\includegraphics[width=\textwidth]{NP_F08_control_chart.png}
\caption{The residual control chart over the monitored record, with alarming
episodes shaded. The line breaks across outages rather than interpolating
through them.}
\label{fig:control-chart}
\end{figure}

\subsection{What this detector can and cannot see}

One large event cannot establish a detector's sensitivity, and the summer-2026
stretch is not used to try. Departures of known size, duration and \emph{shape}
are injected into the uncontaminated residual instead, the charts are re-run, and
an alarm counts only where the uncontaminated run is silent within a
twenty-four-hour response window.

The three shapes correspond to mechanisms a three-leaf stone wall --- two masonry
leaves either side of a weaker rubble-and-mortar core --- can actually produce.
\textbf{Amplitude growth}: the leaves stop acting together, through delamination
at the core interface or loss of through-stones, and the section bends further
under the same daily heating while its timing and mean hold. \textbf{Phase
change}: the thermal path changes rather than the stiffness, as when water
entering the core raises its heat capacity, and the wall answers the same forcing
earlier or later. \textbf{Drift}: creep of the lime mortar under sustained load,
thermal ratcheting of the outer leaf, or foundation settlement --- slow, monotone
and invisible in any single day.

The result is stark, and it is the most operationally important number in this
study. Of the three mechanisms, only amplitude growth is ever detected, and only
at 16~mdeg. A timing shift of up to two hours, which the fitted daily amplitude
of 2.1~mdeg converts to a residual amplitude of 1.07~mdeg, is not detected at any
duration. A drift of up to 20~mdeg per year is not detected at any duration.

\begin{table}[H]
\centering
\caption{Detectability by mechanism, magnitude and persistence, at the stated
false-alarm budget. Read from \texttt{NP\_10\_detectability.csv}.}
\label{tab:detectability}
\begin{tabular}{lrrlr}
\toprule
Mechanism & Magnitude & Duration (h) & Detected & Delay (h) \\
\midrule
\input{../outputs/NP_10_body.tex}\\
\bottomrule
\end{tabular}
\end{table}

Two caveats belong with that table. The amplitude field is not monotone: 16~mdeg
is detected at 6, 24, 72 and 336~hours but not at 168, which is an artefact of
injecting at a single fixed point in the record rather than a property of the
mechanism, and it is reported rather than smoothed. And the whole field is quoted
at a 122-day false-alarm budget; a system willing to tolerate more false alarms
would see smaller departures, and the trade is the operator's to make.

\begin{figure}[H]
\centering
\includegraphics[width=0.8\textwidth]{NP_F09_detectability_amplitude.png}
\caption{Detectability of a growth in the daily swing, the signature of lost
composite action between the leaves.}
\label{fig:detect-amplitude}
\end{figure}

\begin{figure}[H]
\centering
\includegraphics[width=0.8\textwidth]{NP_F09_detectability_phase.png}
\caption{Detectability of a timing shift, the signature of a changed thermal
path such as water in the core. Nothing in the swept range is found.}
\label{fig:detect-phase}
\end{figure}

\begin{figure}[H]
\centering
\includegraphics[width=0.8\textwidth]{NP_F09_detectability_drift.png}
\caption{Detectability of a slow monotone drift, the signature of mortar creep or
settlement. Nothing in the swept range is found.}
\label{fig:detect-drift}
\end{figure}

% ===========================================================================
\section{How far ahead is prediction worth anything?}
\label{sec:forecast}

The target here is the gap-safe one-hour change, formed only between adjacent
accepted observations inside one contiguous run, so that no difference bridges a
gap. After segmentation at a minimum run length of 96~hours, 14,930 hourly rows
in 59 runs remain, evaluated over five expanding whole-segment folds.

Table~\ref{tab:forecast} gives the error of each model and each baseline against
horizon. Every model reaches a mean absolute error between 1.77 and
1.84~mdeg at every horizon from one hour to 48; the zero-change baseline sits at
2.17~mdeg, seasonal-naive at 2.24, and persistence between 2.53 and 3.62.

\begin{table}[H]
\centering
\caption{Forecast error by model and horizon, against three baselines a
monitoring system could run for free. Read from
\texttt{NP\_11\_forecast\_metrics.csv}.}
\label{tab:forecast}
\begin{tabular}{lrrrrrr}
\toprule
Model & Horizon (h) & $n$ & MAE & RMSE & MASE & Coverage \\
\midrule
\input{../outputs/NP_11_body.tex}\\
\bottomrule
\end{tabular}
\end{table}

\begin{figure}[H]
\centering
\includegraphics[width=\textwidth]{NP_F10_skill_vs_horizon.png}
\caption{Forecast error against horizon, by model and baseline.}
\label{fig:skill-horizon}
\end{figure}

An error is not a result until it is compared, and a comparison is not evidence
until it is paired. Table~\ref{tab:skill} gives the fractional reduction in mean
absolute error against each baseline, computed on the timestamps the two share
and with a block bootstrap whose block length is adapted to the horizon, because
these residuals are autocorrelated and an unpaired comparison would report a
significance that is not there.

\begin{table}[H]
\centering
\caption{Paired block-bootstrap skill against each baseline, with 5th and 95th
percentiles of the bootstrap distribution. Read from
\texttt{NP\_12\_skill\_vs\_baseline.csv}.}
\label{tab:skill}
\begin{tabular}{llrrrr}
\toprule
Baseline & Model & Horizon (h) & Skill & $q_{05}$ & $q_{95}$ \\
\midrule
\input{../outputs/NP_12_body.tex}\\
\bottomrule
\end{tabular}
\end{table}

The answer to the third question is therefore that \emph{the horizon at which
forecasting stops beating a trivial baseline was not reached within the range
tested}. Skill against persistence runs from 0.277 at one hour to 0.204 at
48~hours, and against the zero-change baseline from 0.164 to 0.161; every
bootstrap interval excludes zero, including at the longest horizon. The decay
with horizon is mild because the quantity being forecast is a one-hour
increment, not a level: the baselines degrade as fast as the model does.

\subsection{What each predictor actually bought}

A model carrying both autoregressive memory and air temperature cannot say which
of the two earned its skill, which is why the ladder exists. Its result is the
most useful negative finding in this study.

\begin{table}[H]
\centering
\caption{The skill increment each predictor bought over the rung below it. Read
from \texttt{NP\_13\_ablation.csv}.}
\label{tab:ablation}
\begin{tabular}{llrrrr}
\toprule
Over & Model & Horizon (h) & Skill & $q_{05}$ & $q_{95}$ \\
\midrule
\input{../outputs/NP_13_body.tex}\\
\bottomrule
\end{tabular}
\end{table}

\begin{figure}[H]
\centering
\includegraphics[width=\textwidth]{NP_F11_ablation.png}
\caption{The skill increment each predictor bought over the rung below it. Note
the vertical scale: both increments are small fractions of a per cent to a few
per cent, against the 16 to 28\% the models take from the baselines.}
\label{fig:ablation}
\end{figure}

Air temperature, added to autoregression alone, buys 0.0038 of skill at one hour
and less than 0.0013 at every horizon beyond six; at 12 and 48~hours its
bootstrap interval includes zero. Relative humidity, added on top of air
temperature, buys rather more --- 0.023 at one hour, with an interval excluding
zero at every horizon --- though the magnitudes remain small.

The negative control settles the interpretation. Battery voltage added to
autoregression scores 0.273 against persistence at one hour, where autoregression
alone scores 0.274. Adding a channel to this model buys nothing by itself, so the
0.0038 that air temperature buys is a real increment and not an artefact of an
extra input; it is simply a very small one.

This does not contradict Section~\ref{sec:decomposition}. The thermal coupling is
real, was measured here at $-2.56$~\mdegC\ and independently by Study~03 at
$-2.79$, and it explains 12.2\% of the level's variance. It is nonetheless
almost useless for \emph{forecasting the change}, because the response's own
recent history already carries the information the temperature would supply. A
coupling that is real is not obliged to be predictively useful, and the
distinction between the two matters for anyone deciding which sensors a
deployment needs: this record argues for keeping the inclinometer well sampled
rather than for adding environmental channels.

% ===========================================================================
\section{Can the model fill the gaps it was trained around?}
\label{sec:closure}

The operational temptation, once a model predicts a change, is to accumulate its
predictions across a gap and call the result a reconstructed level. That is a
stronger claim than anything measured above: a per-step error that is small and
unbiased still accumulates, and a reconstruction arriving at the wrong level on
the far side of a gap is worse than an admitted absence, because it looks like a
measurement.

The check was to be arithmetical. Where a gap is bracketed by accepted
observations on both sides the true change across it is known, so the predicted
changes can be summed and compared, and a reconstruction is safe only if it
closes.

\textbf{The question is undecidable on this pipeline, and the reason is
structural rather than numerical.} Of the 391 hourly gaps, 390 return
\texttt{incomplete\_estimates} and one \texttt{unbracketed}; none is scorable.
The forecasting model produces predictions only at timestamps that survive
segmentation, and those are by construction the timestamps at which an
observation exists. It therefore never predicts \emph{inside} a gap at all, and
there is nothing to accumulate. Answering the question properly requires rolling
the model forward recursively through each gap, feeding each prediction back as
the next step's history, which this pipeline does not do.

\begin{table}[H]
\centering
\caption{Gap closure summarised by status, which is what decides whether a gap
can be scored at all. None of the 391 gaps is scorable. The full per-gap table
is in \texttt{NP\_14\_gap\_closure.csv}.}
\label{tab:closure}
\begin{tabular}{lrrr}
\toprule
Status & Gaps & Missing slots & Median error \\
\midrule
\input{../outputs/NP_14_body.tex}\\
\bottomrule
\end{tabular}
\end{table}

The honest verdict is therefore that this study does not license gap filling by
accumulation, on the ground that it has not tested it --- which is a weaker
statement than "unsafe" and a much weaker one than "safe". Section~\ref{sec:limitations}
records it as the study's largest piece of unfinished business.

% ===========================================================================
\section{Verdict}
\label{sec:verdict}

\textbf{What is the record made of?} A trend carrying 48.9\% of the variance and
spanning 213~mdeg, a residual carrying 30.4\%, the response to air temperature
12.2\% and the annual cycle 8.4\%. The thermal gain is $-2.56$~\mdegC, which
agrees within 8\% with Study~03's independently measured $-2.79$ and carries the
sign the embankment geometry predicts. That gain survives re-estimation on three
disjoint stretches; the trend does not, disagreeing with itself in sign as well
as magnitude, so the drift is quoted once in-sample at $-2.77$~mdeg/yr and never
extrapolated.

\textbf{Is this reading the expected one?} Yes, to within an expectation whose
mean absolute error is 18.2~mdeg, provided the model is refitted on a schedule.
A frozen fit is not noisy but stale --- 82\% of its error is a single offset ---
and refitting every thirty days removes that. The interval around that
expectation is the weak point: its conformal 90\% band covers 67.7\%, because
the drift the monitoring model deliberately leaves in the residual keeps moving
after the calibration stretch ends. A reading outside that band is evidence; a
reading inside it is weaker evidence than the nominal figure would suggest.

\textbf{How far ahead is prediction worth anything?} Further than this study
tested. Every model beats every baseline out to 48~hours with bootstrap
intervals excluding zero, at a skill of 0.20 against persistence and 0.16
against zero-change. But almost all of that skill is the response's own memory:
air temperature adds 0.4\% at one hour and nothing distinguishable from zero
beyond six, and a battery-voltage control scores the same as autoregression
alone. The environmental channels earn their place in the \emph{decomposition},
not in the \emph{forecast}.

\textbf{And what can the monitor see?} At a measured false-alarm budget of
122~days, only a growth of 16~mdeg in the daily swing. A timing shift of up to
two hours and a drift of up to 20~mdeg per year pass undetected across the whole
range swept. This is the number an operator would need first, and it is
sobering: the configuration tested here detects gross change and not the early,
slow signatures that make monitoring worthwhile.

% ===========================================================================
\section{Limitations}
\label{sec:limitations}

\textbf{The compensation is not independent of the thermal term being fitted.}
The channel modelled here has already been temperature-compensated by Study~01,
and this study fits a thermal response to it. The agreement with Study~03 is
reassuring because that measurement was made on the same compensated channel by
a different method, but neither study can rule out that the compensation and the
fitted term are partly describing the same correction.

\textbf{Two channels were excluded and cover 17\% of the window.} Wall
temperature and solar radiation are the two variables most likely to explain the
residual this study cannot account for, and they were set aside because the
instruments recording them did not exist for three fifths of the window.

\textbf{The battery control is not a silent channel.} Study~03 measured
$r = -0.673$ for supply voltage in the diurnal band, so it marks a conservative
floor rather than a zero. A predictor is credited here only when it clears that
floor, which makes the ablation's negative result stronger, not weaker.

\textbf{Three annual cycles with a 103-day hole cannot identify a yearly term.}
The window spans 3.2 annual cycles, the last of them interrupted for 103~days.
The annual component is retained because leaving a known cause unmodelled would
push it into the residual where a detector would read it as structure, but it is
not well identified, and the burn-in excluded from the reference window in
Section~\ref{sec:monitoring} is a direct symptom of that.

\textbf{The control limit is far outside its textbook range, and rests on two
episodes.} A limit of 14.75 standard deviations is not a conventional control
limit, and it is wide because the residual has a lag-one autocorrelation of
0.995 --- a condition under which EWMA and CUSUM theory does not hold. The
achieved run length of 122~days rests on two excursions and bounds the
false-alarm rate rather than measuring it. Two alternative designs were not
tested and would each be a better answer than a wider limit: charting prewhitened
innovations, and charting a coarser aggregate.

\textbf{Gap filling was not tested at all}, for the structural reason given in
Section~\ref{sec:closure}. Nothing here licenses reconstruction by accumulation.

\textbf{One sensor, one station, one site.} Every number above describes the
inclinometer at station~02 on the wall at Porta di Sant'Ubaldo. The method is
intended to transfer; none of these values is.

\textbf{The summer-2026 stretch is not evidence of detection.} It was excluded
from the reference window as a calibration precaution, and no claim is made
anywhere in this report about whether the detector finds it. A single large event
cannot establish a detection rate, and finding it would show only that the
detector is not broken.

% ===========================================================================
\section{Run metadata}
\label{sec:metadata}

Every parameter and library version behind the artefacts quoted above.

\begin{table}[H]
\centering
\caption{Run metadata. Read from \texttt{NP\_15\_run\_metadata.csv}.}
\label{tab:metadata}
\begin{tabular}{ll}
\toprule
Parameter & Value \\
\midrule
\input{../outputs/NP_15_body.tex}\\
\bottomrule
\end{tabular}
\end{table}

\end{document}
